\chapter{绪论}
学位论文为了反映出作者确已掌握了坚实的基础理论和系统的专门知识，具有开阔的科学视野，对研究方案作了充分论证，因此，有关历史回顾和前人工作的综述分析，以及理论分析等，可以单独成章，用足够的文字叙述。博士学位论文绪论一般不少于 0.8万字，硕士学位论文绪论一般不少于 0.5 万字。\footnote{脚注内容采用宋体小五号字，按两端对齐格式书写，单倍行距，段前段后均空 0 磅。脚注的序号按页编排，不同页的脚注序号不需要连续。详细规定见本页脚注}
\section{研究背景及意义}
正文中某句话需要具体注释、且注释内容与正文内容关系不大时可以采用脚注方式。在正文中需要注释的句子结尾处用 xxxx 样式的数字编排序号，以“上标”字体标示在需要注释的句子末尾。在当页下部书写脚注内容。\footnote{脚注内容采用宋体小五号字，按两端对齐格式书写，单倍行距，段前段后均空 0 磅。脚注的序号按页编排，不同页的脚注序号不需要连续。详细规定见本页脚注}
\section{文献综述}
（1）标注格式：引用参考文献标注方式应全文统一，标注的格式为[序号]，放在引文或转述观点的最后一个句号之前，所引文献序号用小4号Times New Roman体、以上角标形式置于方括号中，如：“…成果\upcite{pritchard1969statistical}”。
（2）序号：标注的序号可选择如下两中方式之一：根据在正文中被首次引用出现的先后次序递增，或者按第一作者姓的英文字母或拼音字母的英文字母顺序递增。如一篇论著在论文中多处引用时，以第一次出现的序号为准；当某一论述同时引证多篇参考文献时，应采用[1，2，…]的格式，标注方括号中的序号按增序排列。
（3）当提及的参考文献为文中直接说明时，其序号应该与正文排齐，如“由文献\cite{pritchard1969statistical}可知：”。

\section{研究内容与研究思路}
\subsection{研究内容}
\subsection{研究方法}
\subsection{技术路线}
\begin{figure}[ht]
\centering
\includegraphics[width=10cm]{cueb.eps}
\caption{The procedures of generating the positive and negative relation images.}
\label{smi}
\end{figure}

\begin{figure}[ht]
\centering
\includegraphics[width=10cm]{cueb.eps}
\caption{The procedures of generating the positive and negative relation images.}
\label{smi}
\end{figure}

\begin{figure}[ht]
\centering
\includegraphics[width=10cm]{cueb.eps}
\caption{The procedures of generating the positive and negative relation images.}
\label{smi}
\end{figure}

\begin{figure}[ht]
\centering
\includegraphics[width=10cm]{cueb.eps}
\caption{The procedures of generating the positive and negative relation images.}
\label{smi}
\end{figure}

\begin{figure}[ht]
\centering
\includegraphics[width=10cm]{cueb.eps}
\caption{The procedures of generating the positive and negative relation images.}
\label{smi}
\end{figure}


\begin{figure}[ht]
\centering
\includegraphics[width=10cm]{cueb.eps}
\caption{The procedures of generating the positive and negative relation images.}
\label{smi}
\end{figure}

\begin{figure}[ht]
\centering
\includegraphics[width=10cm]{cueb.eps}
\caption{The procedures of generating the positive and negative relation images.}
\label{smi}
\end{figure}

\begin{figure}[ht]
\centering
\includegraphics[width=10cm]{cueb.eps}
\caption{The procedures of generating the positive and negative relation images.}
\label{smi}
\end{figure}

\section{论文结构安排}
